\section{Related Work}
\label{sec:related-work}

In this section we review three areas of related work: (i)~treatment effect heterogeneity and external validity, (ii)~decision-focused experiment design, and (iii)~sequential experimentation.

\subsection{External validity and generalizability}

How well experimental findings generalize to broader populations has historically taken a back seat to concerns about internal statistical validity in a number of fields, including the social sciences~\cite{EgamiHa23} and medicine~\cite{AverittWe20}.
Growing evidence suggests this lack of attention to \emph{external validity} is consequential.
Treatment effects often vary substantially across subpopulations, settings, and time periods~\cite{BryanTi21,Tipton21,Meager19,WilliamsCi22}.
The replication crisis in psychology, economics, and biomedical research has revealed that interventions showing promise in initial studies frequently fail to produce similar effects in new contexts~\cite{OpenScienceCollaboration15,CamererDr16,BegleyEl12,CommitteeonReproducibilityandReplicabilityinScienceBo19}.

In resource-constrained settings, many experiments rely on convenience samples.
Driven by limited budgets, tight timelines, and difficulty in recruiting participants from hard-to-reach populations, experimenters select sites based on geographic proximity to their institution, existing relationships, and other factors affecting recruitment~\cite{TiptonFe16,TiptonSp21}.
Enrichment strategies are another common reason for using non-representative samples in clinical trials.
These strategies select a sample from a subpopulation that contains subjects with either an especially positive predicted response to the treatment or an especially poor prognosis, making it easier to detect statistically significant effects~\cite{FDACD19,Thall21,SteingrimssonBe21}.
A third reason for non-representative samples is that certain subpopulations may be more vulnerable or protected, and therefore not eligible for the study.
This is particularly common in biomedical research, where safety is a primary concern~\cite{FDAHH20}.

In response to concerns about external validity, researchers and regulators have advocated for the broader use where possible of representative sampling---constructing experimental samples that mirror the composition of the target population~\cite{Tipton13,FDAHH20}.
This design principle is realized in many large RCTs by attempting to draw a uniform random sample from the target population---for example, in \Cref{sec:case-study} we model the pilot preceding a large education RCT that attempted to draw a random sample of US public schools \cite{YeagerHa19}.
When such random sampling is not feasible, or when small sample sizes might randomly lead to non-representative samples, various methods of nonprobability sampling have been proposed with the goal of making generalizable inferences~\cite{Tipton13,EgamiLe23,Bouyamourn25}.
In contrast to our work, these methods are all focused on a single round of experimentation, and use objectives focused on the external validity of inferences rather than the quality of decision-making.

Another approach to achieving external validity is to attempt to reweight or extrapolate treatment effects from a study sample to a target population~\cite{TiptonOl18,ZhangHu24,LingMo23}.
In this work we model the continuation criterion using the standard approach to significance testing, where inference is conducted for the treatment effect in the study sample, and we leave to future work the design of optimal pilots when scale-up decisions can utilize out-of-sample prediction or extrapolation.

\subsection{Decision-focused experiment design}

While the notion of external validity is concerned with generalization, the question of \emph{what} to generalize is largely independent.
Much of the statistical analysis and design literature is focused on estimating quantities of interest~\cite{Fisher49,Pukelsheim06} or maximizing some measure of information gain~\cite{ChalonerVe95,RainforthFo23}.
In contrast, our work is part of a large body of work, dating back at least to~\textcite{Blackwell51} but especially rich in recent years, that focuses on the value of information in shaping the outcomes of \emph{decisions} instead of the accuracy of estimates, understanding that estimation and decision-making can require significantly different information~\cite{Fernandez-LoriaPr22,Eckles22}.

Much of the decision-focused experimental design literature is focused on Bayesian decision theory, where the design of the experiment and the final decision are jointly optimized over a prior belief distribution~\cite{RyanDr16,NgIm26}.
This approach has been used to derive structural insights into optimal experimentation policies~\cite{RadnerSt84,DeLaraGi07,AzevedoDe20}, and to make policy recommendations for allocating resources to experiments~\cite{FenwickSt20}.
It has also been used to develop specific experiment designs, as in~\textcite{Heath22,FlightJu22,ChenWi13,FrazierPo09,GechterHi24}.
\textcite{GechterHi24} is a work in this vein that, like ours, optimizes the sample composition of an experiment under a prior distribution on heterogeneous treatment effect.
It differs from our work by assuming that adoption decisions are made in a Bayes-optimal manner following the experiment based on the posterior distribution.
Unfortunately, this type of Bayesian decision-making is not feasible in the common scenarios where the experiment designer and decision maker have different information~\cite{BatesJo22,BatesJo23,Min23}, or when frequentist statistical methods are required for reasons of robustness, by scientific norms, or by regulations~\cite{WassersteinLa16,BanerjeeCh20,Recht25}.
Our work differs from standard Bayes-optimal in that it takes into account the widespread use of frequentist significance testing as an evaluation criterion, while still allowing for optimizing the experiment design using prior information.
We ask:~\emph{how do current statistical testing norms affect optimal experiment design?}

Unlike the Bayesian approach which requires a prior distribution, minimax-regret experiment design instead considers worst-case decision performance against some set of possible outcomes, and has similarly been used to derive optimal site selection or sample composition designs~\cite{OleaPr25,HuZh24}.
Similarly to our work, \textcite{HuZh24} studies how to allocate sampling effort across subpopulations in a randomized experiment to optimize downstream welfare under heterogeneous effects.
In contrast to our approach, they present a solution under a minimax-regret criterion, and for a single experiment stage where adoption may differ across groups. Their design does not depend qualitatively on the budget, which we believe reflects the fact that the minimax-optimal adoption criterion they model does not require statistical significance.

\subsection{Sequential experimentation}

Like decision-focused experiment design, the design of \emph{sequential} experiments has been studied for a long time~\cite{Wald47}, but has recently received special focus as computational methods have made their optimization more feasible~\cite{CheNa23,ChapelleLi11,TrellaZh22} and domains like online platforms have created more opportunities for real-world adaptive experimentation~\cite{AzevedoDe20,XuDu18}.
Such sequential or adaptive experiments are promising in their ability to improve over static designs in terms of the tradeoff between experiment size or duration and quality, as well as by controlling risk during the experiment and other objectives.

We focus on what we believe is a first-order effect in many multi-stage experiments: \emph{the screening effect of early-stage experiments on what interventions receive further study}.
This screening effect has been studied in the context of designing phase II / efficacy trials in drug development, with various works optimizing some combination of sample sizes for phase II and phase III, the decision rule to continue to phase III, and adaptive stopping and sample sizes within stage II~\cite{RossellMu07,DingRo08,KirchnerKi16,ErdmannKi20}.
We are not aware of any work optimizing the sample composition, i.e.\ the distribution of covariates, in a screening trial under a decision- or utility-focused objective.

While adaptivity is helpful in many settings, in many domains it is not clear how or even if the design of a follow-up RCT depends on the outcomes of a previous experiment.
For example, efficiency trials in medical and education research often must use representative or otherwise constrained designs, and are often run by different researchers and institutions than pilot studies, limiting the range of possible adaptivity~\cite{FDA19,DepartmentofEducationBo22}.

As we describe in \Cref{sec:Model}, we fix all the decision rules downstream of the pilot study including the continuation and adoption criteria (statistical significance tests) and the follow-up design (a fixed sampling design independent of the pilot outcome, for example a uniform random sample with a given size).
While we model a multi-stage experimentation process, our optimization problem only involves a single decision---the pilot design.
